\documentclass[10pt,conference]{IEEEtran}

\usepackage{amsmath}
\usepackage{booktabs}
\usepackage{graphicx}
\usepackage{multirow}
\usepackage{url}
\usepackage{xcolor}
\usepackage{enumitem}

\newcommand{\todo}[1]{\textcolor{red}{[TODO: #1]}}
\newcommand{\tool}{VRAgent~2.0}
\newcommand{\lc}{LC}
\newcommand{\mc}{MC}
\newcommand{\coigo}{CoIGO}

\begin{document}

\title{Coverage-Guided Multi-Agent Exploration for Testing Unity XR Applications}

\author{Anonymous Authors}

\maketitle

\begin{abstract}
Extended reality (XR) applications combine spatial interaction, scene hierarchies,
physics objects, Unity events, and runtime-generated state. These properties make
automated test generation difficult: a syntactically valid action may target an
unreachable object, miss a required component, invoke an unavailable method, or
fail semantically because a hidden precondition has not been satisfied. Recent
large language model (LLM)-based approaches can generate interaction plans from
project context, but one-shot generation provides little opportunity to repair
invalid plans or to adapt exploration based on runtime feedback.

This paper presents \tool, a retrieval-augmented, contract-based multi-agent
framework for closed-loop testing of Unity XR applications. \tool decomposes XR
test generation into specialized agents: a Planner proposes action units, a
Verifier checks structural and semantic validity, an Executor deterministically
executes actions through an online Unity bridge, and an Observer interprets
runtime traces, coverage signals, and failure evidence. A shared blackboard and a
Gate Graph record discovered states, blocked interactions, inferred conditions,
and scheduler priorities. The system uses project retrieval over Unity scenes,
GameObject hierarchies, component metadata, and C\# scripts to construct
targeted context for each agent. To support online decision making before final
Unity Code Coverage reports are available, \tool also maintains a runtime
novelty proxy based on new successful action patterns, events, method calls, and
state-changing objects.

We plan to evaluate \tool on a suite of Unity XR applications by comparing it
against one-shot VRAgent generation, no-retrieval generation, random and greedy
exploration, and ablated variants of the multi-agent loop. The primary metrics
are line coverage, method coverage, coverage of interesting game objects,
successful interaction rate, gate resolution rate, redundancy, and token/runtime
cost. \todo{Fill in quantitative results after running the complete benchmark.}
\end{abstract}

\begin{IEEEkeywords}
XR testing, Unity testing, large language models, multi-agent systems,
coverage-guided testing, retrieval-augmented generation
\end{IEEEkeywords}

\section{Introduction}

XR applications are increasingly used in training, education, healthcare,
entertainment, and industrial simulation. Many of these applications are built
with Unity and the XR Interaction Toolkit, where behavior is distributed across
GameObjects, Prefabs, physics components, UnityEvents, and C\# scripts. Testing
such applications requires more than invoking UI callbacks. A test system must
move through a 3D scene, grab objects, place them into spatial targets, trigger
buttons and sockets, open panels, satisfy interaction preconditions, and observe
whether scripts actually executed.

Automating this process is challenging for at least three reasons. First, object
references in Unity are brittle: scene objects, Prefab instances, and runtime
objects require stable resolution mechanisms. Second, XR interactions are often
gated by hidden state, such as inventory contents, door locks, active UI panels,
or puzzle flags. Third, a technically executable action may not increase code
coverage if it repeats an already tested interaction or targets an object with no
meaningful behavior.

Our prior VRAgent workflow addressed part of this problem by extracting Unity
scene information and asking an LLM to generate JSON test plans that could be
imported and executed by a Unity-side agent. However, that workflow is largely
one-shot: a plan is generated from static context, then executed. If the plan
contains invalid references, misses preconditions, or produces no new coverage,
the system lacks a principled mechanism to diagnose the failure and adapt the
next action.

This paper argues that effective XR testing requires a closed-loop architecture.
Rather than treating LLMs as one-shot plan generators, \tool treats them as
participants in a contract-based pipeline guided by static project knowledge,
runtime execution traces, failure hypotheses, and coverage feedback. The system
iteratively selects objects, retrieves local context, generates action units,
checks them, executes them in Unity, observes the outcome, and updates a shared
world state for the next round.

The paper makes the following contributions:

\begin{itemize}[leftmargin=*]
  \item We characterize common failure modes in LLM-based XR application testing,
  including invalid object references, missing XR components, invalid method
  calls, semantic no-effect actions, and unsatisfied interaction gates.
  \item We present a retrieval-augmented multi-agent architecture that separates
  planning, verification, execution, and observation through explicit schemas and
  a shared blackboard.
  \item We introduce a Gate Graph and failure-to-condition reasoning mechanism
  for converting blocked XR interactions into actionable exploration goals.
  \item We implement an online Unity execution bridge and a runtime novelty proxy
  that provide immediate feedback for coverage-guided exploration.
  \item We design an empirical evaluation over multiple Unity XR applications,
  with baselines and ablations that isolate the effect of retrieval, verification,
  observation, gate reasoning, revisiting, and coverage feedback.
\end{itemize}

\todo{After experiments, replace the final contribution with concrete numerical
findings, e.g., percentage coverage improvement and cost overhead.}

\section{Background and Motivation}

\subsection{Unity XR Application Structure}

Unity XR applications are organized as scenes containing GameObjects. A
GameObject may have physics components, colliders, XR interactables, scripts,
UnityEvents, children, tags, and layers. Behaviors are often distributed: a
button object may invoke a method on a door controller; a key may update an
inventory manager; a panel may only become active after a puzzle flag changes.

This structure creates a mismatch between object-level testing and behavior-level
coverage. Testing a GameObject once is insufficient if the behavior requires a
sequence of interactions across multiple objects. Conversely, repeatedly
interacting with an object after its methods have already been covered wastes the
testing budget.

\subsection{VRAgent 1.0}

VRAgent 1.0 uses FileID-based object resolution to stabilize scene references.
It imports a JSON test plan consisting of task units and action units. Supported
action types include Grab, Trigger, Transform, Move, and Socket. At runtime,
Unity-side components such as XRGrabbable, XRTriggerable, and XRTransformable are
dynamically attached to target objects to execute actions.

The limitation is that VRAgent 1.0 executes a generated plan rather than
exploring interactively. When an action fails or produces no new behavior, the
feedback is not used to repair the plan or choose a better next object.

\subsection{Motivating Failure Example}

Consider an escape-room scene with a locked door and a key inside a drawer. A
one-shot LLM plan may attempt to trigger the door first. The action is
structurally valid because the door object and method exist, but it has no
semantic effect because the player has not collected the key. A useful testing
system should infer that the door is blocked by a missing condition, prioritize
the drawer/key objects, and revisit the door after the condition is satisfied.

\todo{Replace this toy example with a real case study from one subject app.}

\section{Approach}

Figure~\ref{fig:overview} gives an overview of \tool. The pipeline is organized
around a shared blackboard. Each iteration consumes the current exploration goal
and previous observer output, selects a target object, retrieves relevant project
context, generates a candidate plan, verifies it, executes it in Unity, observes
the outcome, and updates the blackboard and Gate Graph.

\begin{figure*}[t]
\centering
\fbox{\begin{minipage}{0.94\textwidth}
\centering
\vspace{0.8em}
\textbf{Placeholder for Figure 1: System Overview}\par
\vspace{0.5em}
Unity Project Analysis $\rightarrow$ Retrieval Layer $\rightarrow$ Blackboard\\
Scheduler $\rightarrow$ Planner $\rightarrow$ Static/Semantic Verifier
$\rightarrow$ Executor $\rightarrow$ Unity Bridge $\rightarrow$ Unity Runtime\\
Observer $\rightarrow$ Coverage, Failure Hypotheses, GateGraph, Scheduler Bias
\vspace{0.8em}
\end{minipage}}
\caption{Overview of \tool. Replace this placeholder with an architecture diagram
showing Python agents, retrieval, blackboard, Gate Graph, and Unity runtime.}
\label{fig:overview}
\end{figure*}

\subsection{Project Retrieval Layer}

The retrieval layer indexes Unity scene graphs, GameObject hierarchies,
component metadata, script sources, tag/layer logic, and special APIs such as
\texttt{GameObject.Find} and \texttt{Instantiate}. Given a target object and a
goal, it constructs a context pack for the Planner containing object summaries,
nearby interactables, relevant scripts, special logic, gate hints, and recent
failures. For the Verifier, it constructs structured caches for object existence,
component information, and method names.

The key design principle is that agents should receive local, goal-relevant
context rather than the entire project. This reduces prompt size and makes the
evidence used by each agent explicit.

\subsection{Contract-Based Agents}

Each agent communicates through a schema. The Planner returns candidate action
units and an expected reward. The Verifier returns an executable score, errors,
and patched actions. The Executor returns traces, exceptions, and coverage
delta. The Observer returns bug signals, state deltas, failure hypotheses, gate
hints, and strategy recommendations.

This separation gives each component a narrow responsibility. The Planner is not
trusted to validate its own output; the Executor does not generate or repair
actions; the Observer does not execute tests. This design supports ablation and
debugging because each stage can be inspected independently.

\subsection{Verification and Repair}

The static Verifier checks structural correctness: valid FileIDs, required
fields, component requirements, method availability, duplicate actions, and
repeated invalid targets. If errors are found, it returns structured repair
signals to the Planner. A semantic verifier optionally uses an LLM to critique
whether the plan makes sense given known scene facts and gate chains.

\subsection{Online Execution}

The Executor communicates with Unity through a TCP bridge. Before executing an
action, it imports the required FileIDs into Unity. Unity executes the action and
returns a trace containing pre-state, post-state, events, exceptions, duration,
and a success flag. This enables Python-side agents to observe behavior without
waiting for a full batch plan to finish.

\subsection{Observation and Weak Oracle}

The Observer acts as a weak oracle. It detects exceptions and console errors,
interprets state changes, distinguishes completed trigger actions from true
no-effect failures, and infers possible missing conditions from trace and log
patterns. For example, evidence containing terms such as ``locked'', ``key'',
``inactive'', or ``not found'' is mapped to structured failure types such as
\texttt{locked}, \texttt{need\_item}, \texttt{ui\_hidden}, and
\texttt{reference\_invalid}.

\subsection{Gate Graph}

The Gate Graph records successful and failed transitions. A failed edge is a
candidate gate. When the Observer infers a condition for a failed edge, the
Scheduler and Planner can focus on satisfying that condition rather than
continuing a linear object traversal.

\begin{figure}[t]
\centering
\fbox{\begin{minipage}{0.44\textwidth}
\centering
\vspace{0.8em}
\textbf{Placeholder for Figure 2: Gate Graph}\par
\vspace{0.5em}
$S_0 \xrightarrow{Trigger(Door)} blocked:locked$\\
$S_0 \xrightarrow{Grab(Key)} S_1$\\
$S_1 \xrightarrow{Trigger(Door)} S_2$\\
\vspace{0.8em}
\end{minipage}}
\caption{Example Gate Graph for a locked-door interaction chain. Replace with a
real graph extracted from an experiment.}
\label{fig:ggraph}
\end{figure}

\subsection{Coverage-Guided Scheduling}

The exploration controller uses a three-mode state machine: Expand, Exploit,
and Recover. Expand searches for new objects or UI branches. Exploit focuses on
inferred gate conditions. Recover is triggered after repeated zero-novelty
iterations.

In addition to final Unity Code Coverage reports, the Executor maintains a
runtime novelty proxy:
\begin{equation}
r_t = 0.4 \Delta LC_p + 0.3 \Delta MC_p + 0.3 \Delta CoIGO_p,
\end{equation}
where $\Delta LC_p$ counts new successful action patterns and events,
$\Delta MC_p$ counts new method call units, and $\Delta CoIGO_p$ counts newly
changed interesting objects. This proxy is used for online scheduling; final
reported results should use Unity Code Coverage measurements.

\section{Implementation}

\tool is implemented as a Python package and a Unity runtime package. The Python
package contains the controller, agents, retrieval layer, scheduler, Gate Graph,
and bridge client. The Unity package contains an online executor that receives
commands, resolves FileIDs, dynamically attaches XR components, executes actions,
captures object states, and drains console logs.

Table~\ref{tab:modules} summarizes the implementation modules.

\begin{table*}[t]
\centering
\caption{Main implementation modules.}
\label{tab:modules}
\begin{tabular}{lll}
\toprule
Component & Location & Responsibility \\
\midrule
Controller & \texttt{vragent2/controller.py} & Closed-loop orchestration \\
Retrieval & \texttt{vragent2/retrieval/} & Scene/script/context retrieval \\
Planner & \texttt{vragent2/agents/planner.py} & Candidate action generation \\
Verifier & \texttt{vragent2/agents/verifier.py} & Static and semantic validation \\
Executor & \texttt{vragent2/agents/executor.py} & Deterministic execution and traces \\
Observer & \texttt{vragent2/agents/observer.py} & Weak oracle and strategy recommendation \\
Scheduler & \texttt{vragent2/scheduling/object\_scheduler.py} & Target object selection \\
Gate Graph & \texttt{vragent2/graph/gate\_graph.py} & Gate and state transition memory \\
Unity bridge & \texttt{vragent2/bridge/unity\_bridge.py} & TCP protocol client \\
Unity runtime & \texttt{VRAgentOnline.cs} & In-editor action execution \\
\bottomrule
\end{tabular}
\end{table*}

\todo{Add implementation size: Python LOC, C\# LOC, number of commands, number of
schemas, and supported action types.}

\section{Evaluation Plan}

\subsection{Research Questions}

\begin{description}[leftmargin=*]
  \item[RQ1] Does \tool improve code and interaction coverage compared with
  one-shot LLM test generation and non-LLM exploration baselines?
  \item[RQ2] Which components contribute most to effectiveness?
  \item[RQ3] Does closed-loop observation reduce invalid, redundant, or
  no-effect actions?
  \item[RQ4] What is the runtime and token cost of multi-agent exploration?
  \item[RQ5] How does \tool behave on gated interaction chains in XR scenes?
\end{description}

\subsection{Subject Applications}

We will evaluate \tool on Unity XR applications spanning room-scale navigation,
object manipulation, escape-room puzzles, and simulation tasks. Candidate
subjects include BowlingVR, EE-Room, EscapeTheRoom, maze, Parkinson-VR,
UnityCityView, VGuns, VR-Basics, VR-Room, vr-firefighter-simulator, and
wheelchair-sim.

\begin{table*}[t]
\centering
\caption{Subject applications. Fill values after running the analysis pipeline.}
\label{tab:subjects}
\begin{tabular}{lrrrrr}
\toprule
Subject & Scenes & GameObjects & MonoBehaviours & UnityEvents & LOC \\
\midrule
\todo{Subject A} & -- & -- & -- & -- & -- \\
\todo{Subject B} & -- & -- & -- & -- & -- \\
\todo{Subject C} & -- & -- & -- & -- & -- \\
\todo{Subject D} & -- & -- & -- & -- & -- \\
\bottomrule
\end{tabular}
\end{table*}

\subsection{Baselines}

We compare against the following baselines:

\begin{itemize}[leftmargin=*]
  \item \textbf{VRAgent 1.0 one-shot}: the previous one-shot plan generation
  workflow.
  \item \textbf{NoTODG}: generation without the target object dependency graph.
  \item \textbf{Random}: randomly chooses objects and action patterns.
  \item \textbf{Greedy static}: prioritizes objects by static script/component
  features without runtime feedback.
  \item \textbf{LLM-only online}: uses an LLM online but removes verification,
  observer feedback, and Gate Graph memory.
\end{itemize}

\todo{Add any SOTA fuzzing/RL/game-testing baseline that can be fairly run on the
same Unity XR subjects. If a direct baseline is unavailable, justify why and use
closest runnable alternatives.}

\subsection{Metrics}

The primary metrics are line coverage (LC), method coverage (MC), and coverage
of interesting game objects (CoIGO). Secondary metrics include action success
rate, semantic success rate, gate solved count, duplicate action rate, invalid
target rate, no-effect action rate, token cost, and wall-clock runtime.

\subsection{Ablations}

\begin{table*}[t]
\centering
\caption{Planned ablation variants.}
\label{tab:ablations}
\begin{tabular}{lcccccc}
\toprule
Variant & RAG & Verifier & Observer & GateGraph & Revisit & Runtime proxy \\
\midrule
Full \tool & yes & yes & yes & yes & yes & yes \\
w/o RAG & no & yes & yes & yes & yes & yes \\
w/o Verifier & yes & no & yes & yes & yes & yes \\
w/o Observer & yes & yes & no & no & no & yes \\
w/o GateGraph & yes & yes & yes & no & yes & yes \\
w/o Revisit & yes & yes & yes & yes & no & yes \\
w/o Runtime Proxy & yes & yes & yes & yes & yes & no \\
\bottomrule
\end{tabular}
\end{table*}

\subsection{Expected Result Tables}

\begin{table*}[t]
\centering
\caption{Overall effectiveness. Replace placeholders with mean and standard
deviation over repeated runs.}
\label{tab:results}
\begin{tabular}{lrrrrrr}
\toprule
Technique & LC & MC & CoIGO & Success & Gates & Cost \\
\midrule
Random & -- & -- & -- & -- & -- & -- \\
Greedy static & -- & -- & -- & -- & -- & -- \\
NoTODG & -- & -- & -- & -- & -- & -- \\
VRAgent 1.0 & -- & -- & -- & -- & -- & -- \\
\tool & -- & -- & -- & -- & -- & -- \\
\bottomrule
\end{tabular}
\end{table*}

\section{Discussion}

\subsection{Why Closed-Loop Feedback Matters}

One-shot LLM plans fail because they cannot observe whether an interaction was
useful. The runtime trace gives the system evidence about which actions succeed,
which objects change, which methods are invoked, and which conditions are
missing. The Observer converts this evidence into planner instructions and
scheduler bias, enabling the system to revisit important objects after new
preconditions are inferred.

\subsection{Runtime Proxy vs. Ground-Truth Coverage}

The runtime novelty proxy is not a replacement for Unity Code Coverage. It is an
online reward used to guide exploration before offline coverage reports are
available. In the final evaluation, all coverage claims should be based on Unity
Code Coverage summaries. The proxy should be analyzed separately as a guidance
signal.

\subsection{Limitations}

\tool currently relies on static extraction quality and FileID resolution. Some
runtime-generated objects may still be missed. The Observer uses heuristic
failure-to-condition inference and may misclassify failures. The LLM agents can
produce inconsistent action semantics, although the Verifier mitigates many
structural errors. Finally, the current runtime proxy approximates coverage and
must be validated against actual Unity coverage data.

\section{Threats to Validity}

\textbf{Internal validity.} Results may depend on prompt templates, LLM model
versions, Unity version, and scene setup. We will use fixed seeds where possible,
log all prompts and responses, and repeat runs to reduce nondeterminism.

\textbf{Construct validity.} LC and MC measure code execution but may not fully
capture XR interaction quality. We therefore include CoIGO, gate resolution, and
semantic success metrics.

\textbf{External validity.} The subject applications may not represent all XR
genres or industrial XR systems. We mitigate this by selecting diverse projects
with different interaction patterns.

\textbf{Conclusion validity.} Statistical comparisons require repeated runs and
appropriate tests. \todo{Add statistical method: Wilcoxon signed-rank, effect
size, confidence intervals, or bootstrap depending on final data.}

\section{Related Work}

\todo{Add and cite related work in the following areas: automated GUI testing,
game testing, Unity/VR testing, search-based testing, coverage-guided fuzzing,
LLM-based test generation, retrieval-augmented generation for SE, and multi-agent
LLM systems. Use real bibliographic entries before submission.}

\section{Conclusion}

This paper presented \tool, a closed-loop, retrieval-augmented multi-agent
framework for testing Unity XR applications. By combining structured project
retrieval, contract-based planning and verification, online Unity execution,
weak-oracle observation, Gate Graph memory, and runtime coverage guidance, \tool
addresses key limitations of one-shot LLM test generation. The planned empirical
evaluation will quantify its impact on code coverage, interaction coverage,
failure reduction, and exploration cost across multiple XR applications.

\section*{Artifact Availability}

\todo{For double-anonymous submission, provide an anonymized artifact link or a
statement explaining artifact availability. Include setup instructions, subject
apps, scripts, and raw logs. Remove identifying paths and names.}

\section*{AI Assistance Disclosure}

\todo{Before submission, add the required ACM/IEEE generative-AI disclosure in a
way consistent with ICSE 2027 double-anonymous policy.}

\bibliographystyle{IEEEtran}
\begin{thebibliography}{1}

\bibitem{todo-vr-testing}
\todo{Add canonical references for VR/XR application testing.}

\bibitem{todo-llm-testing}
\todo{Add references for LLM-based software test generation.}

\bibitem{todo-coverage-guided}
\todo{Add references for coverage-guided fuzzing and search-based testing.}

\bibitem{todo-gui-testing}
\todo{Add references for GUI/mobile/game testing.}

\end{thebibliography}

\end{document}
